\documentclass[11pt, letterpaper]{article}
\input{../../homework.tex}
\usepackage{titlepageBU}
\title{Problem Set 8}
\courseID{PY 355}
\courseSection{A1}
\begin{document}
\makeproblem
\section*{Problem 1}
\begin{enumerate}
	\item 
		\[
			\int_{-2}^{1} x^3\delta (x+1) \,\mathrm{d} x=(-1)^3=-1
		.\]
	\item 
		\[
			\int_{-\infty }^{y} \delta (x) \,\mathrm{d} x=
			\begin{cases}
				0,y<0\\
				1,y\geq 0
			\end{cases}
		.\]
		\item
			\[
				\int_{0}^{3\pi /2} \cos x\delta (x-\pi ) \,\mathrm{d} x=\cos \pi =-1
			.\]
		\item 
			\[
				\int_{\mathbb{R} ^3} \left( r^2+\vec{r} \cdot \vec{a} +a^2 \right) \delta ^{(3)} \left( \vec{r} -\vec{a}  \right)  \,\mathrm{d} r=3a^2
			.\]
\end{enumerate}
\section*{Problem 2}
\begin{enumerate}
	\item Let $u=a(x-x_0)$. Then $du=adx$, so
		\[
			\int_{-\infty }^{\infty } f(x)\delta (a(x-x_0)) \,\mathrm{d} x=\int_{-\infty }^{\infty }\frac{1}{a} f\left( x_0+\frac{u}{a}  \right) \delta (u) \,\mathrm{d} u=f\left( x_0+\frac{0}{a}  \right) \frac{1}{a} 
		.\]
		Note that the above equations are mild abuse of notation; if $a<0$, then the bounds in $u$ are the reverse of that in $x$, and thus shifting them back induces an extra factor of $-1$. Thus, if $a<0$, then the bounds change, inducing an extra negative. We can thus write the more proper solution as
		\[
			\int_{-\infty }^{\infty } f(x)\delta (a(x-x_0)) \,\mathrm{d} x=\frac{f(x_0)}{\left| a \right| } 
		.\]
	\item We have
		\[
			\int_{-\infty }^{\infty } f(x)\delta (a(x-x_0)) \,\mathrm{d} x=\int_{-\infty }^{\infty } f(x)\frac{1}{\left| a \right| }\delta (x-x_0) \,\mathrm{d} x
		.\]
		In order for this equality to hold for all integrable functions $f$, we must have
		\[
			\frac{1}{\left| a \right| } \delta (x-x_0)=\delta (a(x-x_0))
		.\]
	\item We thus have
		\[
			\delta (-x)=\delta (-1(x-0))=\frac{1}{\left| -1 \right| } \delta (x-0)=\delta (x)
		.\]
		Therefore, $\delta $ is even.
\end{enumerate}

\end{document}
